\documentclass[12pt]{article}
\usepackage[utf8]{inputenc}
\usepackage{amsmath, amssymb}
\usepackage[spanish]{babel}
\usepackage{geometry}
\geometry{a4paper, margin=2cm}

% ==========================================
% COMANDOS DE ESTRUCTURA PARA EL PARSER (AI)
% ==========================================
% Estos comandos permiten que la IA lea automáticamente tu documento LaTeX
% y lo convierta directamente a los archivos JSON del sistema, soportando
% cualquier jerarquía (preguntas sueltas, anidadas con contexto, anidadas sin contexto).

% 1. Configuración general de la Evaluación
% Uso: \evalMeta{ID}{Título}{Descripción}{Curso}{Universidad}{Año}{Semestre}{Nombre de la Evaluación}{Minutos}
% Opciones para "Nombre de la Evaluación": Control, Prueba, Certamen, Examen
\newcommand{\evalMeta}[9]{
  \begin{center}
    {\Large \textbf{#2}}\\[0.2cm]
    \textbf{Curso:} #4 \quad \textbf{Universidad:} #5 \quad \textbf{Año:} #6 - #7\\[0.1cm]
    \textbf{Evaluación:} #8 \quad \textbf{Tiempo límite:} #9 minutos\\[0.3cm]
    \textit{#3}
  \end{center}
  \vspace{0.5cm}
}

% 2. Etiquetas generales del ensayo (separadas por coma)
\newcommand{\evalTags}[1]{
  \textbf{Etiquetas Generales:} #1 \vspace{0.5cm}
}

% 3. Bloque para agrupar subpreguntas bajo un mismo contexto general
% Uso: \begin{bloqueContexto}{Aquí va el enunciado general compartido}
\newenvironment{bloqueContexto}[1]{
  \vspace{0.5cm}\noindent\hrulefill\vspace{0.2cm}\\
  \noindent \textbf{Contexto General:}\\
  #1
  \vspace{0.3cm}
}{
  \vspace{0.5cm}
}

% 4. Entorno de Pregunta (aplica tanto para sueltas como para subpreguntas)
% Uso: \begin{pregunta}{ID_Pregunta}{Puntaje}{Dificultad}{Temas}
\newenvironment{pregunta}[4]{
  \vspace{0.3cm}\noindent\textbf{Pregunta (#2 pts) - Dificultad: #3}\\
  \textbf{Temas:} #4\\[0.2cm]
}{
  \vspace{0.3cm}
}

% 5. Sub-elementos de la pregunta (Pistas, Soluciones, Pauta y Errores)

% Usa \numeracion{} para indicar el número de la pregunta en la UI (ej: 1. a) i)
\newcommand{\numeracion}[1]{
  \vspace{-0.3cm}\noindent\textbf{Numeración:} #1 \\
}

\newcommand{\pista}[1]{
  \noindent \textbf{\textsl{Pista:}} #1 \\
}

\newenvironment{solucion}{
  \vspace{0.2cm}
  \noindent \textbf{Solución Paso a Paso:}\\[0.1cm]
}{
  \vspace{0.2cm}
}

\newcommand{\criterio}[2]{
  \noindent $\square$ \textbf{Criterio (#2 pts):} #1 \\
}

\newcommand{\errorfrecuente}[1]{
  \noindent \textbf{Error Frecuente:} #1 \\
}


% ==========================================
% INICIO DEL DOCUMENTO
% ==========================================
\begin{document}

\evalMeta
  {examen-ejemplo}
  {Examen de Ejemplo (Estructuras Jerárquicas)}
  {Demostración de cómo armar preguntas con y sin nidos.}
  {Cálculo}
  {Universidad}
  {2026}
  {Otoño}
  {Examen}
  {90}

\evalTags{Derivadas, Integrales}

% ---------------------------------------------------------
% CASO 1: PREGUNTAS ANIDADAS CON ENUNCIADO COMÚN
% ---------------------------------------------------------
\begin{bloqueContexto}{
Respecto de la función
$$ f(x,y) = \begin{cases} \frac{x^2y}{x^2+y^2} & (x,y) \neq (0,0) \\ 0 & (x,y) = (0,0) \end{cases} $$
}

    \begin{pregunta}{q1-a}{5}{Media}{Derivadas parciales}
    \numeracion{1. a)}
    Calcule las derivadas parciales.
    \pista{Usa regla de la cadena.}
    \end{pregunta}

    \begin{pregunta}{q1-b}{5}{Media}{Diferenciabilidad}
    \numeracion{1. b)}
    Decida si la función es diferenciable en el origen.
    \end{pregunta}

\end{bloqueContexto}


% ---------------------------------------------------------
% CASO 2: PREGUNTAS ANIDADAS PERO INDEPENDIENTES (SIN CONTEXTO COMÚN)
% ---------------------------------------------------------
% Aquí simplemente pones las preguntas seguidas y usas \numeracion para agruparlas visualmente en la plataforma web, sin necesidad de un bloqueContexto.

\begin{pregunta}{q2-a}{10}{Alta}{Integrales}
\numeracion{2. a)}
Considere la integral triple... (enunciado específico de a)
\end{pregunta}

\begin{pregunta}{q2-b}{10}{Alta}{Polares}
\numeracion{2. b)}
Calcule la integral doble... (enunciado específico de b, totalmente ajeno a la 2a)
\end{pregunta}


% ---------------------------------------------------------
% CASO 3: PREGUNTA SUELTA (SIN ANIDAR)
% ---------------------------------------------------------
\begin{pregunta}{q3}{15}{Media}{Regla de la Cadena}
\numeracion{3.}
Suponga que $f,g$ son diferenciables. Muestre que se cumple la identidad...
\end{pregunta}


\end{document}
